Figure~\ref{fig:atlas-convex-chord} shows the chord inequality that characterizes convexity.

\begin{figure}[!htbp]
\centering
\begin{tikzpicture}[>=Stealth]
\begin{axis}[width=11cm,height=5.8cm,axis lines=middle,ticklabel style={font=\small},label style={font=\small},legend style={font=\scriptsize,draw=black!20,fill=white},xlabel={$x$},ylabel={$f(x)$},xmin=-1.5,xmax=2.4,ymin=-0.3,ymax=4.7,xtick={-1,0.5,2},ytick={1,2.5,4}]
\addplot[figblue,very thick,domain=-1.3:2.1,samples=120] {x^2};
\addplot[figred,thick] coordinates {(-1,1)(2,4)};
\draw[figgreen,thick,<->] (axis cs:0.5,0.25) -- (axis cs:0.5,2.5);
\addplot[only marks,figblue,mark=*] coordinates {(-1,1)(2,4)(0.5,0.25)};
\addplot[only marks,figred,mark=*] coordinates {(0.5,2.5)};
\end{axis}
\end{tikzpicture}
\caption{For the convex function $f(x)=x^2$, the line segment joining two graph points lies above the graph. At the midpoint of $x_1=-1$ and $x_2=2$, the function value is $1/4$ whereas the mean of the endpoint values is $5/2$. The vertical gap illustrates the convexity inequality.}
\label{fig:atlas-convex-chord}
\end{figure}